\scanpage{97}
% archive foliation: 96
Quelles connexions \unsure{affines} obtient-on ainsi ?
\deleted{Choisissons (localt) une \unsure{uniformisante} $t$ de $X$, d'où une base $dt$ de $\omega\simeq\Omega$, et un splitting de $P^1_{X/S}(\Omega)\simeq P^1_{X/S}(\omega)\simeq P^1_{X/S}(\mathcal{O}_X)$ (par le choix de la base $dt$), d'où une connexion \inserted{projective} correspondante.}
On peut supposer $\mathcal{H}$ splitté localement, avec base $\delta$ de $\Omega$, et une connexion de référence $c_0$ décrite par
\[
\begin{aligned}
c_0(\underset{\substack{\Vert\\\lambda\delta}}{\omega_1})&=\delta\,d\lambda+\omega_1(-2\bar\omega_0),&\bar\omega_0&\in\Gamma\Omega^{-1}\otimes\omega^{\unsure{1}},\\
c_0(\underset{\substack{\Vert\\\mu\cdot1}}{\omega_0})&=1\,d\mu+\omega_0(\bar\omega_0),\\
c_0(\underset{\substack{\Vert\\\unsure{\gamma|\nu}\delta^{-1}}}{\omega_{-1}})&=\delta^{-1}d\unsure{\gamma|\nu}\,\deleted{\unsure{\omega}}.
\end{aligned}
\]
Considérons une section $s$ de $\mathcal{L}$ s'envoyant sur section $1_X$ de $\mathcal{O}_X$
\[
s=s_0+\omega_1\qquad(s_0=(0,1))
\]
soit
\[
\left\{\begin{array}{l}
\deleted{\unsure{\psi_s}:\mathcal{O}_X\longrightarrow\mathcal{L}}\\
\nu_s:\Omega^{-1}\longrightarrow\mathcal{H}
\end{array}\right.
\qquad\text{i.e. }\nu_s\in\Gamma(\mathcal{H}\otimes\Omega)
\]
\margin{\deleted{\unsure{$\psi_s\in\Gamma(\mathcal{L})$}}}
\deleted{l'élément}\inserted{l'hom} correspondant. \deleted{\unsure{et}} On a
\[
\deleted{s=\exp\Theta_{\omega_1}\cdot s_0=s_0+[-\omega_1,s_0]+\unsure{[-\omega_1]}}
\]
\[
\nu_s=\exp\Theta_{\omega_1}\nu_{s_0}
=\nu_{s_0}+\underbrace{[-\omega_1,\nu_{s_0}]}_{\unsure{-2\omega_1\nu_{s_0}}}
+\underbrace{\frac12[-\omega_1,[-\omega_1,\nu_{s_0}]]}_{\unsure{+\omega_1^2\nu_{s_0}}}
\]
\margin{\unsure{pour $s_0=s$}}
\deleted{d'où la question est si $c(\nu_s)$ est} de calculer la composante \inserted{$\bar\omega_2(s)=\varphi(\omega_1)$} de $c(s)$, suivant la \deleted{base}\inserted{$\Omega\otimes\omega$} \deleted{\unsure{que la}} décomposition
\[
\mathcal{H}\simeq\Omega+\deleted{\unsure{\exp\Theta}}\mathcal{O}_X\cdot s+\nu_s(\Omega^{-1})
\]
i.e.
\[
(\exp\Theta_{\omega_1})\otimes\operatorname{id}_{\omega}\bigl(c(s)\bigr).
\]
\deleted{et la forme cherchée (pour corriger) est}
\[
\deleted{\unsure{\bar\omega_2(s)=-\varphi(\omega_1)\in\Omega\otimes\omega}.}
\]
\[
\deleted{\underbrace{c(\nu_s)=\bar\omega_2(s)+\unsure{\bar\omega_1(s)}\,\unsure{s\otimes\bar\omega_1(1)}+(\nu_s\otimes\operatorname{id})(\bar\omega_{-1})}_{\substack{\Omega\otimes\omega\qquad\omega\qquad\Omega^{-1}\otimes\omega\\=\text{composante de }\exp\Theta_{+\omega_1}(c(s))}}}
\]
\margin{\deleted{$\unsure{\operatorname{id}\otimes c(s_0)}$}\inserted{$\unsure{-\omega_1}$}}
Sauf erreur, on trouve
\[
s=s_0+\omega_1\longmapsto c_0+\bigl(\delta\,d(\unsure{\omega_1\delta^{-1}})+\omega_1^2\bigr).
\]

% right half, continuing the left half
Prenons $\delta$ de façon que son image dans $\omega$ soit $dt$, $t$ une uniformisante, on trouve sauf erreur de calcul --
\[
s_0+\lambda\,dt\longmapsto c_0^{\deleted{\unsure{dt}}}+\left(\frac{d\lambda}{dt}+\lambda^2\right)(dt)^2
\]
\margin{facteur $2$ !}
Donc quand une conn. proj. sur $X/S$ est définie par (\unsure{au-delà} de la \deleted{forme}\inserted{conn.} proj. de réf. $c_0$) par une $2$-forme $\varphi(dt)^2$, la question si on peut trouver une section horizontale de $P$ disjointe de la section donnée revient à celle de trouver des solutions $\lambda$ de l'équation
\[
\frac{d\lambda}{dt}+\lambda^2=\varphi
\]
(équation diff. non lin. du $1^{\mathrm{er}}$ ordre). On sait, si $X$ courbe lisse sur corps $k$ de car nulle, qu'il y a au plus une solution ayant une valeur et une dérivée données en un pt; \unsure{s'il y en a une infinité}, mais transcendantes (!), en \uline{algébrique} on peut \unsure{avoir} --

Les pôles de $\lambda$ \inserted{\unsure{nous}} correspondent à des pts de rencontre de $\sigma$ \inserted{section} (défini par $\lambda$) et de $\sigma'$ --

\uline{1.6.} Intrinsèquement, le $P$ canonique est $\simeq\mathbb{P}(P^1_{X/S}(\underline{\omega}))$, où $\underline{\omega}=\underline{\Omega}^1_{X/S}$. Donc une section de $P$ s'identifie à un couple
\[
(L,\gamma)
\]
$L$ Module inversible sur $X$, et
\[
\gamma:P^1_{X/S}(\omega)\longrightarrow L\qquad\text{épi. lin.}
\]
Or un hom. \deleted{\unsure{lin.}}\inserted{\unsure{l'op.}} $\gamma:P^1_{X/S}(\omega)\longrightarrow L$ équivaut à
